% =====================================================================
\section{The whole picture on one page}
\fig[0.8\textwidth]{fig_overview.png}{The working hypothesis as a single chain. Grey arrows are cause and effect; orange arrows are feedback.}
The hypothesis the brochure is built around has six steps (fig.~1).
\begin{enumerate}
\item The cortex contains a small group of inhibitory cells, the \emph{PV cells}. They hold the excitation of the rest in check and set the fast rhythm of the network.
\item PV cells work at their limit: many inputs, high firing rate, a large energy turnover. A by-product of that turnover is a flux of oxidants.
\item Enzymes protect the cell from oxidants. Part of that defence is produced only while the cell's NMDA receptor --- its activity sensor --- is working properly.
\item If the receptor works weakly, through heredity or through the environment, there is less defence and there are more oxidants. Oxidants in turn suppress the receptor: the loop closes on itself.
\item An overloaded PV cell does not die; it \emph{falls silent}, producing less of its working proteins. Inhibition weakens, the cortex is disinhibited, and the load on the surviving cells rises: a second loop.
\item An over-excited hippocampus, through a chain of intermediate structures, raises the responsivity of the dopamine system --- the bridge to the positive symptoms.
\end{enumerate}
\begin{note}
Every arrow in this chain has been shown separately in experiment: in cell culture, in brain slices or in animals. The chain as a whole has never been. A single hypothesis is set out here --- weak NMDA-receptor function together with impaired oxidative defence of PV cells. Its competitors are not reviewed.
\end{note}
\begin{ineq}
A circle that amplifies itself is positive feedback. If its gain is large enough, the system acquires two stable states and a threshold between them, and the transition between them is abrupt rather than gradual. That is what makes a sudden onset possible at all.
\end{ineq}
